\section{Conclusion}
\label{sec:conclusion}

AI is showing up in the macro data, but not where the productivity narrative expects. Measured output per hour in exposed sectors has not accelerated, and total employment in exposed industries has not changed. What has changed is the composition of hiring: exposed sectors employ markedly fewer young workers relative to prime-age workers, a pattern that began with COVID and, in our preferred specification, steepened after ChatGPT. The 1995--2004 IT boom shows the opposite sign, which is itself evidence that AI is a different kind of technology for white-collar work. Whether the gain is reaching shareholders is not yet visible; the labor share of exposed sectors has returned to about its 2019 level after a COVID-era spike.

For an investor, the result is an early labor-market signal that has not yet become a cash-flow signal, and that is drowned out by rate exposure at sector level. A long-short strategy on the junior-substitution signal did not earn excess returns from 1996 to 2026. We would not deploy it as a stand-alone strategy. We would use the signal in two ways: as a filter when choosing among AI-exposed equities, and as an early warning. A clear fall in the labor share of exposed sectors over 2026--27 would turn the signal into a cash-flow story worth trading; a reversal of the junior decline would argue that the 2022--25 pattern was cyclical after all.

Four extensions would sharpen the analysis: event windows around BEA and QWI release dates using daily returns; employment-weighted exposure; firm-level labor and margin data, which would replace 16 sector portfolios with hundreds of stocks; and a markup measure that separates pricing power from substitution in the labor share.
